\item \textbf{Problema de repaso: Energía potencial electrostática de un núcleo.}
Considere un modelo del núcleo atómico donde la carga positiva neta $Ze$ se encuentra uniformemente distribuida por toda una esfera de radio $R$.
\begin{enumerate}
    \item Demuestre, integrando la densidad de energía electrostática $u = \frac{1}{2} \varepsilon_0 E^2$ en todo el espacio infinito, que la energía potencial eléctrica total requerida para ensamblar el núcleo puede escribirse como:
    $$ U = \frac{3 k_e Z^2 e^2}{5 R} $$
    \item Muestre paso a paso el cálculo de las integrales para la región interior ($r < R$) y para la región exterior ($r > R$).
\end{enumerate}